\chapter{3D Rendering: Models and Meshes}
\label{ch:models-meshes}

\section{The runtime hierarchy}

\cnaclass{Model} owns a \cnaclass{ModelBoneCollection} and a \cnaclass{ModelMeshCollection};
each \cnaclass{ModelMesh} owns a \cnaclass{ModelMeshPartCollection} and a
\cnaclass{ModelEffectCollection}; each \cnaclass{ModelMeshPart} references a shared
\cnaclass{VertexBuffer} / \cnaclass{IndexBuffer} pair and an \cnaclass{Effect}. This is the
same shape as real XNA, and \cnaclass{Model}'s \emph{runtime} API --- as opposed to how a
model actually gets loaded, \S\ref{sec:model-loading} below --- is fully audited against FNA
with zero known open gaps as of the most recent fix (a \texttt{rootBoneIndex} correction).
\cnaclass{Model} exposes \texttt{CopyAbsoluteBoneTransformsTo},
\texttt{CopyBoneTransformsFrom} / \texttt{To}, and a convenience \texttt{Draw(world, view,
projection)}; one small, intentionally-kept deviation from FNA is that CNA's
\texttt{Copy*BoneTransforms*} methods loop by \texttt{Bones.Count} rather than trusting the
caller's array length the way FNA does --- strictly safer against an oversized destination
array, never a compatibility problem for correctly-sized callers.

\cnaclass{ModelBone} carries \texttt{Name}, \texttt{Index}, a get/set \texttt{Transform},
its \texttt{Parent}, and a \texttt{Children} collection built via a \texttt{NOXNA}
\texttt{AddChild}. \cnaclass{ModelMesh} carries a get/\emph{set} \texttt{BoundingSphere}
(the setter was added specifically to support the real \texttt{.xnb} content pipeline
described below) and a get/\emph{set} \texttt{ParentBone} --- both properties that only
matter once something actually populates them, which is exactly where the two different
model loaders in this codebase diverge sharply.

\section{A worked example: two ways to draw a loaded model}

The simplest path is the convenience method already named above:

\begin{lstlisting}
model.Draw(worldMatrix, camera.View(), camera.Projection());
\end{lstlisting}

This is correct for a model with no bone hierarchy to speak of, but a real, multi-bone model
(a character, a vehicle with independently-animated wheels) needs each mesh drawn with its
\emph{own} absolute bone transform composed into world space first --- the classic pattern,
grounded directly in the real \texttt{CopyAbsoluteBoneTransformsTo} and
\texttt{IEffectMatrices} signatures:

\begin{lstlisting}
std::vector<Matrix> boneTransforms(model.getBonesProperty().getCountProperty());
model.CopyAbsoluteBoneTransformsTo(boneTransforms);

for (ModelMesh& mesh : model.getMeshesProperty()) {
    Matrix meshWorld = boneTransforms[mesh.getParentBoneProperty()->getIndexProperty()]
                      * worldMatrix;

    for (Effect* effect : mesh.getEffectsProperty()) {
        if (auto* matrices = dynamic_cast<IEffectMatrices*>(effect)) {
            matrices->setWorldProperty(meshWorld);
            matrices->setViewProperty(camera.View());
            matrices->setProjectionProperty(camera.Projection());
        }
    }
    mesh.Draw(); // issues the actual GPU draw calls per ModelMeshPart -- Chapter 9 covers
                 // the underlying GraphicsDevice::DrawIndexedPrimitives path this reaches.
}
\end{lstlisting}

\texttt{boneTransforms[mesh.getParentBoneProperty()->getIndexProperty()]} is the one line
worth pausing on: it looks up the mesh's \emph{own} parent bone's absolute transform, not bone
0 or the model's own root --- which is exactly why the earlier ``two loaders'' distinction
matters practically here, not just architecturally. Every mesh drawn through a real
\texttt{.xnb}-loaded model gets its correct, distinct \texttt{ParentBone}; every mesh loaded
through the older \texttt{.model.json} path silently gets the same default bone regardless of
what the source data says, so this exact loop --- copied verbatim from a real XNA tutorial ---
would render every part of a \texttt{.model.json}-loaded model rigidly welded together, with no
error or warning that anything was wrong.

\section{Two loaders, two very different stories}
\label{sec:model-loading}

This is the single most important thing to understand about \cnaclass{Model} in CNA, and it
is easy to miss if you only read the runtime class headers: \textbf{there are two separate,
independent \cnaclass{Model} loaders}, and \texttt{ContentManager::Load<Model>()} always
tries the newer, real one first.

\begin{description}
  \item[The real \texttt{.xnb} \texttt{ModelReader}] (\texttt{CNA::Internal::Xnb::ModelReader})
        is wire-compatible with real XNA/MonoGame/FNA compiled model assets. It reconstructs
        the full bone hierarchy, assigns each mesh's real \texttt{ParentBone}, computes a real
        \texttt{BoundingSphere}, and correctly resolves shared \cnaclass{VertexBuffer}/%
        \cnaclass{IndexBuffer} / \cnaclass{Effect} resources across meshes that share them ---
        everything below in this section that describes a gap does \emph{not} apply to this
        loader.
  \item[The older \texttt{.model.json} loose-file loader] (\texttt{ModelTypeReader}) is what
        every gap below describes. It always synthesizes exactly one \cnaclass{ModelBone}
        (index~0), even when the source JSON's own \texttt{"bones"} array lists more --- only
        the first entry's name is ever read, and no \texttt{AddChild} hierarchy is ever built.
        Every mesh's \texttt{ParentBone} is left at its default regardless of what the JSON
        says, even though the runtime constructor has supported an explicit parent-bone
        argument since the fix mentioned above --- the loader simply never calls it.
        \texttt{BoundingSphere} is never computed, staying at a degenerate zero radius, and
        neither \texttt{Model.Tag} nor any \texttt{ModelMesh.Tag} is ever set. Resource
        sharing is also absent: every mesh gets a freshly-allocated
        \cnaclass{VertexBuffer} / \cnaclass{IndexBuffer}, and every part using anything other
        than the plain \texttt{"BasicEffect"} string gets its own freshly-constructed
        \cnaclass{Effect} rather than a shared one --- there is no dependency-graph-aware
        resource deduplication of the kind FNA's own \texttt{ReadSharedResource<T>}
        performs. At the time of writing, this loader has \emph{zero} test coverage of its
        own, confirmed by a repository-wide search.
\end{description}

Whether this distinction matters to you depends entirely on which loader your assets go
through --- a real \texttt{.xnb} asset pipeline (this book's migration chapter) sidesteps
every gap in this section; a hand-authored \texttt{.model.json} asset inherits all of them.

\section{Skinning: two independent systems, not one}

CNA ships two separate, unrelated systems for animated meshes, and conflating them is an
easy mistake to make. \cnaclass{AnimationPlayer} (\texttt{NOXNA}, mirroring Microsoft's own
unshipped XNA ``Skinned Model Sample'' reference code rather than anything in the framework
assembly itself) drives a bone-track timeline: \texttt{StartClip(clip)},
\texttt{Update(TimeSpan, relativeToCurrentTime, loop=true)}, and
\texttt{GetBoneTransforms()} / \texttt{GetWorldTransforms()} / \texttt{GetSkinTransforms()} ---
the last of which is exactly what feeds \cnaclass{SkinnedEffect::SetBoneTransforms}
(Chapter~\ref{ch:stock-effects}). Its skinning data (bone count, hierarchy, bind pose,
inverse bind pose, animation clips) is stashed on a loaded \cnaclass{Model}'s own
\texttt{Tag} --- the same convention the original Microsoft sample used, since real XNA's
\cnaclass{Model} has no dedicated skinning-data property of its own to hold it in.

\cnaclass{MorphTargetEXT} (also \texttt{NOXNA}) is a completely independent, glTF-style
morph-target blending system: CPU-side re-blend plus a full vertex-buffer re-upload on every
call to \texttt{SetMorphWeightsEXT}, a deliberate simplicity-over-throughput tradeoff.
Because glTF's ``weights'' animation channel targets a mesh instance node rather than a bone
index, this system has no relationship to \cnaclass{AnimationPlayer}'s bone timeline at all;
it supports linear, step, and real cubic-spline (Hermite) interpolation, and attaches itself
via a loaded \cnaclass{ModelMeshPart}'s own \texttt{Tag} --- the same attachment convention,
applied to a different object in the hierarchy, for a different purpose.

A third, still more separate system exists for the Avatar subsystem specifically:
\texttt{SkinnedModelEXT} (its own \texttt{.skeleton.bin} / \texttt{.clip.bin} binary formats) is
deliberately not built on \cnaclass{Model} / \cnaclass{ModelBone} / \cnaclass{ModelMesh} at all,
has its own full test coverage, and is covered alongside the rest of the Avatar system in
this book.

\subsection{A worked example: driving AnimationPlayer, adapted from its own test suite}

\texttt{AnimationPlayerTests.cpp} builds its fixtures around the smallest rig that can
actually exercise bone-hierarchy composition: two bones, where bone~1 is a child of bone~0
offset by $(0,1,0)$ in bind pose, and bone~0's own track moves it from the origin to
$(2,0,0)$ over one second. Adapted directly from that fixture:

\begin{lstlisting}
SkinningData data;
data.BoneCount = 2;
data.SkeletonHierarchy = {-1, 0};                              // bone 0 = root, bone 1's parent = 0
data.BindPose = {Matrix::getIdentityProperty(),
                  Matrix::CreateTranslation(Vector3(0, 1, 0))};
data.InverseBindPose = {Matrix::getIdentityProperty(), Matrix::getIdentityProperty()};

BoneTrackEXT track;
track.BoneIndex = 0;
track.Keys.push_back(Keyframe{System::TimeSpan::FromSeconds(0.0), Vector3(0, 0, 0)});
track.Keys.push_back(Keyframe{System::TimeSpan::FromSeconds(1.0), Vector3(2, 0, 0)});

AnimationClip clip;
clip.Duration = System::TimeSpan::FromSeconds(1.0);
clip.Tracks.push_back(track);
data.AnimationClips["Move"] = clip;

AnimationPlayer player(data);
player.StartClip(data.AnimationClips["Move"]);

// Every frame:
player.Update(gameTime.getElapsedGameTimeProperty(), /*relativeToCurrentTime=*/true);
skinnedEffect.SetBoneTransforms(player.GetSkinTransforms());
\end{lstlisting}

Two details this fixture makes concrete rather than abstract: \texttt{SkeletonHierarchy} is a
flat parent-index array (\texttt{-1} means ``no parent, this is a root''), the same
representation glTF and most DCC export pipelines use, not a pointer-based tree; and
\texttt{GetWorldTransforms()} versus \texttt{GetSkinTransforms()} answer two different
questions --- the former is each bone's absolute transform (what \S\ref{sec:model-loading}'s
manual mesh-draw loop above would want for a non-skinned mesh attached to a bone), the latter
is pre-multiplied by each bone's own \texttt{InverseBindPose}, which is specifically the form
\cnaclass{SkinnedEffect::SetBoneTransforms} expects, since GPU skinning needs the bone's
\emph{delta} from bind pose, not its absolute pose.

\section{Collections, uniformly}

\cnaclass{ModelBoneCollection}, \cnaclass{ModelMeshCollection},
\cnaclass{ModelMeshPartCollection}, and \cnaclass{ModelEffectCollection} all follow the same
shape: integer indexing via \texttt{operator[]}, name-based indexing where it makes sense,
\texttt{Count}, \texttt{Contains}, \texttt{TryGetValue}, and standard iterators --- the same
collection idiom already seen on \cnaclass{GameComponentCollection}
(Chapter~\ref{ch:game-loop}) and \cnaclass{CurveKeyCollection}
(Chapter~\ref{ch:math-core-types}), applied consistently a third time here.
